\section{习题课 \#1（9月13日）}

\subsection{Gauss--Jordan 消元与线性方程组}

本次课首先讨论含参数的非齐次线性方程组。核心工具是将方程组写成增广矩阵，并通过可逆的初等行变换化为行阶梯形或简化行阶梯形。

\begin{example}
考虑
\begin{equation}
\begin{cases}
 x_1-x_2+x_3=1,\\
 x_1-x_2-x_3=3,\\
 2x_1-2x_2-x_3=5.
\end{cases}
\end{equation}
其增广矩阵为
\begin{equation}
\left(
\begin{array}{ccc|c}
1&-1&1&1\\
1&-1&-1&3\\
2&-2&-1&5
\end{array}
\right).
\end{equation}
经初等行变换可化为
\begin{equation}
\left(
\begin{array}{ccc|c}
1&-1&0&2\\
0&0&1&-1\\
0&0&0&0
\end{array}
\right).
\end{equation}
因此解的一般形式为
\begin{equation}
 x_1=x_2+2,
 \qquad x_3=-1,
 \qquad x_2\in\F.
\end{equation}
\end{example}

\begin{definition}
矩阵的三类初等行变换为
\begin{enumerate}
  \item 交换两行；
  \item 将某一行乘以非零常数；
  \item 将某一行的常数倍加到另一行。
\end{enumerate}
它们都可逆，因此不会改变线性方程组的解集。
\end{definition}

\begin{remark}
消元过程中每一步都必须是等价变形。对于不等式，这一点尤其容易被忽略。例如
\begin{equation}
 \frac{2x+3}{x+4}\leq 1
 \quad\Longleftrightarrow\quad
 \frac{x-1}{x+4}\leq 0,
\end{equation}
不能在不知道 $x+4$ 符号的情况下直接同乘 $x+4$。正确答案为
\begin{equation}
 x\in(-4,1].
\end{equation}
同理，在求方程 $f(x)=0$ 的解集时，若某一步只保留了单向蕴含，就必须额外检查增根或漏根。Gauss--Jordan 消元之所以可靠，正是因为每一步初等行变换均可逆。
\end{remark}

\subsection{解的三种可能与数域的影响}

线性方程组的解集只有三种可能：无解、唯一解或无穷多解。行化简后可按如下方式判断。
\begin{enumerate}
  \item 若出现形如 $(0,\ldots,0\mid c)$ 且 $c\neq0$ 的行，则方程组无解。
  \item 若每个未知量所在列都有主元，则方程组至多有一个解；在相容时即为唯一解。
  \item 若方程组相容且存在自由变量，则有无穷多解。
\end{enumerate}

所取数域也会影响结论。例如同一个整数系数方程组，在 $\Q$、$\R$ 与有限域 $\F_p$ 上可能具有不同的秩和不同的解数。行变换中的除法只能除以当前数域中的非零元素。

\subsection{典型计算}

\begin{exercise}[线性规划的一个简单例子]
设三种原料的某种成分含量分别为 $12\%$、$15\%$ 和 $22\%$。若总量固定，并且第三种原料的用量固定为 $5$，则约束可整理为
\begin{equation}
 x_1+x_2=5,
 \qquad x_3=5,
 \qquad x_1,x_2,x_3\geq0.
\end{equation}
目标函数为
\begin{equation}
 L=0.12x_1+0.15x_2+0.22x_3.
\end{equation}
求 $L$ 的可能范围。
\end{exercise}

\begin{solution}
由 $x_1=5-x_2$ 和 $0\leq x_2\leq5$，
\begin{equation}
 L=0.12(5-x_2)+0.15x_2+0.22\cdot5
   =1.7+0.03x_2.
\end{equation}
故
\begin{equation}
 1.7\leq L\leq1.85.
\end{equation}
\end{solution}

\begin{exercise}
设 $n\geq2$，$A$ 是元素均为 $\pm1$ 的 $n\times n$ 整数矩阵。证明 $\det A$ 为偶数。
\end{exercise}

\begin{proof}
模 $2$ 意义下有 $1\equiv-1$，故 $A$ 的所有元素均同余于 $1$。于是 $A$ 在 $\F_2$ 上的各行相同，秩至多为 $1<n$，从而
\begin{equation}
 \det A\equiv0\pmod2.
\end{equation}
故 $\det A$ 为偶数。
\end{proof}

\begin{exercise}[插值方程的相容性]
问是否存在二次多项式
\begin{equation}
 p(x)=ax^2+bx+c
\end{equation}
同时通过四个给定点。将插值条件写成线性方程组后，可用增广矩阵判断相容性。
\end{exercise}

\begin{solution}
以原稿中的点
\begin{equation}
 (0,2),\quad (-4,1),\quad (-1,3),\quad (1,2)
\end{equation}
为例，代入后得到
\begin{equation}
\begin{cases}
 c=2,\\
 16a-4b+c=1,\\
 a-b+c=3,\\
 a+b+c=2.
\end{cases}
\end{equation}
由后两式及 $c=2$ 得
\begin{equation}
 a-b=1,
 \qquad
 a+b=0,
\end{equation}
从而 $a=\tfrac12$、$b=-\tfrac12$。但此时
\begin{equation}
 16a-4b+c=8+2+2=12\neq1,
\end{equation}
故方程组不相容，不存在满足条件的二次多项式。这个例子说明，当插值条件的数目超过自由系数数目时，必须额外检查数据之间的相容性。
\end{solution}

\begin{exercise}[一类对称耦合方程组]
设 $a_i\neq0$，且
\begin{equation}
 1+\sum_{i=1}^n a_i^{-1}\neq0.
\end{equation}
求解
\begin{equation}
 (1+a_i)x_i+\sum_{j\neq i}x_j=b_i,
 \qquad i=1,\ldots,n.
\end{equation}
\end{exercise}

\begin{solution}
令
\begin{equation}
 X_0=\sum_{j=1}^n x_j.
\end{equation}
第 $i$ 个方程化为
\begin{equation}
 a_i x_i+X_0=b_i,
 \qquad
 x_i=a_i^{-1}(b_i-X_0).
\end{equation}
求和得
\begin{equation}
 X_0=\sum_{i=1}^n a_i^{-1}b_i-X_0\sum_{i=1}^n a_i^{-1},
\end{equation}
故
\begin{equation}
 X_0=\frac{\sum_{i=1}^n a_i^{-1}b_i}
 {1+\sum_{i=1}^n a_i^{-1}}.
\end{equation}
因此
\begin{equation}
 x_i=a_i^{-1}b_i-
 \frac{a_i^{-1}\sum_{j=1}^n a_j^{-1}b_j}
 {1+\sum_{j=1}^n a_j^{-1}},
 \qquad i=1,\ldots,n.
\end{equation}
\end{solution}

\begin{exercise}[含参数非齐次方程组]
讨论
\begin{equation}
\begin{cases}
 x_1+x_2+x_3+x_4+x_5=1,\\
 3x_1+2x_2+x_3+x_4-3x_5=a,\\
 x_2+2x_3+2x_4+6x_5=3,\\
 5x_1+4x_2+3x_3+3x_4-x_5=b
\end{cases}
\end{equation}
的相容性。
\end{exercise}

\begin{solution}
对增广矩阵作消元，最后两条独立的相容性条件为
\begin{equation}
 a=0,
 \qquad b=2.
\end{equation}
在此条件下有三个自由变量，可取 $x_3,x_4,x_5$，并得到
\begin{equation}
 x_1=x_3+x_4+5x_5-2,
 \qquad
 x_2=-2x_3-2x_4-6x_5+3.
\end{equation}
若 $(a,b)\neq(0,2)$，消元中出现矛盾行，方程组无解。
\end{solution}

\begin{exercise}[含参数齐次方程组]
讨论
\begin{equation}
\begin{cases}
 x_1+x_2+x_3=0,\\
 2x_1+x_2-a x_3=0,\\
 x_1-2x_2-3x_3=0
\end{cases}
\end{equation}
的解。
\end{exercise}

\begin{solution}
系数矩阵行列式为
\begin{equation}
 \det\begin{pmatrix}
 1&1&1\\
 2&1&-a\\
 1&-2&-3
 \end{pmatrix}
 =-3a-2.
\end{equation}
因此当 $a\neq-\frac23$ 时只有零解；当 $a=-\frac23$ 时秩下降，存在一维非零解空间。
\end{solution}

\begin{exercise}[两个行列式恒等式]
设 $a_i,b_i,c_i\in\F$。证明
\begin{equation}
 \det\begin{pmatrix}
 a_1-b_1&b_1-c_1&c_1-a_1\\
 a_2-b_2&b_2-c_2&c_2-a_2\\
 a_3-b_3&b_3-c_3&c_3-a_3
 \end{pmatrix}=0,
\end{equation}
并证明
\begin{equation}
 \det\begin{pmatrix}
 a_1+b_1&b_1+c_1&c_1+a_1\\
 a_2+b_2&b_2+c_2&c_2+a_2\\
 a_3+b_3&b_3+c_3&c_3+a_3
 \end{pmatrix}
 =2\det\begin{pmatrix}
 a_1&b_1&c_1\\
 a_2&b_2&c_2\\
 a_3&b_3&c_3
 \end{pmatrix}.
\end{equation}
\end{exercise}

\begin{proof}
第一个矩阵三列之和为零，故列向量线性相关。对第二个矩阵，记右端三列为 $a,b,c$，则左端三列为 $a+b,b+c,c+a$。相应的列变换矩阵为
\begin{equation}
 T=\begin{pmatrix}1&0&1\\1&1&0\\0&1&1\end{pmatrix},
 \qquad \det T=2,
\end{equation}
所以左端行列式等于 $2\det(a,b,c)$。
\end{proof}

\begin{exercise}[由两个已知解反推参数]
已知
\begin{equation}
 (0,1,0)^T,
 \qquad (-3,2,2)^T
\end{equation}
均为方程组
\begin{equation}
\begin{cases}
 x_1-x_2+2x_3=-1,\\
 3x_1+x_2+4x_3=1,\\
 ax_1+bx_2+cx_3=d
\end{cases}
\end{equation}
的解。求使该方程组有无穷多解的参数条件。
\end{exercise}

\begin{solution}
代入第一个解得 $b=d$；代入第二个解得
\begin{equation}
 -3a+2b+2c=d.
\end{equation}
结合 $b=d$ 可得
\begin{equation}
 3a=b+2c.
\end{equation}
前两行系数向量线性无关，因此它们给出秩为 $2$ 的相容系统。第三行要使秩不增加，必须是前两行的线性组合，这恰好等价于
\begin{equation}
 b=d,
 \qquad 3a=b+2c.
\end{equation}
在此条件下未知数个数为 $3$ 而秩为 $2$，故解集为一条仿射直线。
\end{solution}
